% ══════════════════════════════════════════════════════════════════════════════
%  SECTION – Model Extensions
% ══════════════════════════════════════════════════════════════════════════════

\section{Model Extensions}
\label{sec:extensions}

The base model can be extended by activating any subset of five independent
additions: \textbf{M} (transport modes), \textbf{R} (RTI types),
\textbf{H} (hub nodes), \textbf{P} (explicit product flows), and \textbf{L}
(linearisation of the queue-time products).
{Extensions M, R, H, and L are fully self-contained and
can be activated in any combination without restriction.
Extension P, however, is \emph{not} independent: it requires Extension R to
be active, since product-level flows $NF_{ijpr}$ are indexed over RTI
types $r \in R$, which are introduced exclusively by Extension R.
The dependency structure is depicted in Figure~\ref{fig:ext-deps}.}

% ── Dependency figure ─────────────────────────────────────────────────────────
\begin{figure}[h]
\centering
\begin{tikzpicture}[
  extbox/.style={
    rectangle, rounded corners=6pt,
    draw=blue!60!black, fill=blue!8,
    minimum width=2.4cm, minimum height=1.1cm,
    align=center, font=\small\bfseries
  },
  depbox/.style={
    rectangle, rounded corners=6pt,
    draw=orange!80!black, fill=orange!12,
    minimum width=2.4cm, minimum height=1.1cm,
    align=center, font=\small\bfseries
  },
  reqarrow/.style={
    -{Stealth[length=7pt]}, line width=1.5pt, color=red!70!black
  }
]
  %% Independent extensions (top row)
  \node[extbox] (M) at (0, 0)   {M\\\scriptsize Transport Modes};
  \node[extbox] (R) at (3, 0)   {R\\\scriptsize RTI Types};
  \node[extbox] (H) at (6, 0)   {H\\\scriptsize Hub Nodes};
  \node[extbox] (L) at (9, 0)   {L\\\scriptsize Linearisation};

  %% Dependent extension (below R)
  \node[depbox] (P) at (3, -2.6) {P\\\scriptsize Product Flows};

  %% Dependency arrow
  \draw[reqarrow] (R.south) -- (P.north)
    node[midway, right=5pt, font=\scriptsize\itshape, color=red!70!black]
      {requires};

  %% Background shading for the independent group
  \begin{scope}[on background layer]
    \node[draw=gray!40, fill=gray!5, rounded corners=8pt, dashed,
      fit=(M)(R)(H)(L), inner sep=12pt,
      label={[font=\scriptsize\itshape, color=gray!60]above:%
             independent -- can be freely combined}] {};
  \end{scope}

  %% Legend
  \begin{scope}[shift={(-0.3,-4.4)}]
    \node[extbox, minimum width=1.6cm, minimum height=0.65cm,
          font=\scriptsize] (l1) at (0,0) {};
    \node[right=0.2cm of l1, font=\scriptsize] {Independent extension};
    \node[depbox, minimum width=1.6cm, minimum height=0.65cm,
          font=\scriptsize] (l2) at (6,0) {};
    \node[right=0.2cm of l2, font=\scriptsize] {Requires another extension};
  \end{scope}
\end{tikzpicture}
\caption{Extension dependency graph.  Extensions \textbf{M}, \textbf{R},
\textbf{H}, and \textbf{L} are mutually independent and may be combined
freely.  Extension \textbf{P} (explicit product flows) requires Extension
\textbf{R} (multiple RTI types) to be active, because product-level demand
$d_{ijp}$ is converted into RTI quantities via the type-specific factor
$\eta_{pr}$.}
\label{fig:ext-deps}
\end{figure}

Extensions M, R, H, and P each introduce a new set with its index, replacing
the root variables of the model by adding {an} additional
index. Extension~L is qualitatively different: it does not add an index to
existing variables but instead discretises the continuous queue time $q_{ij}$
over a finite set $\mathcal{Q}_{ij}$ of admissible integer values in order to
{discretise the model.}
The result is a pure MILP formulation with no nonlinear terms.

Each extension operates through two mechanisms.

\begin{enumerate}
  \item \textbf{Index addition or variable introduction.}
        Table~\ref{tab:index-map} lists, for every symbol, which extensions
        add an index to the variable or parameter or introduce a new symbol
        (\textbf{New}), either a variable or parameter.
        Symbols marked $\circ$ in a column are unaffected; symbols marked
        $\times$ are replaced by derived expressions
        (Section~\ref{sec:ext-new-constraints}).

  \item \textbf{Constraint transformation.}
        Table~\ref{tab:constraint-transform} records, for every base
        constraint, what each active extension does to it. Extensions not
        listed for a constraint leave it unchanged ($\circ$). Unless
        specified, a constraint is transformed automatically if a variable
        that appears in it is transformed, by iterating over the new set,
        that is, if an index is added to a variable or parameter.
\end{enumerate}

% ─────────────────────────────────────────────────────────────────────────────
%  TABLE 1 — INDEX MAP
% ─────────────────────────────────────────────────────────────────────────────

\begin{longtable}{llccccc}
\caption{Index promotion and variable introduction map.
  $+m$/$+r$: extension adds the indicated index to this symbol.
  \textbf{New}: symbol is introduced by this extension.
  $\circ$: unaffected.
  $\times$: replaced by a derived expression.
  For Ext.~L the $k$ subscript ranges over $\mathcal{Q}_{ij}$; if Ext.~M is
  also active, replace $(i,j)$ by $(i,j,m)$ throughout.}
\label{tab:index-map}\\
\toprule
\textbf{Symbol} & \textbf{Base form}
  & \textbf{M} & \textbf{R} & \textbf{H} & \textbf{P} & \textbf{L} \\
\midrule
\endfirsthead
\multicolumn{7}{l}{\small\itshape (Table~\ref{tab:index-map} continued)}\\
\toprule
\textbf{Symbol} & \textbf{Base form}
  & \textbf{M} & \textbf{R} & \textbf{H} & \textbf{P} & \textbf{L} \\
\midrule
\endhead
\midrule
\multicolumn{7}{r}{\small\itshape Continued on next page}\\
\endfoot
\bottomrule
\endlastfoot

\multicolumn{7}{l}{\textit{Sets}} \\
$\mathcal{M}$
  & ---
  & \textbf{New} & $\circ$ & $\circ$ & $\circ$ & $\circ$ \\
$R,\;R_{ij}$
  & ---
  & $\circ$ & \textbf{New} & $\circ$ & $\circ$ & $\circ$ \\
$H\subseteq V$
  & ---
  & $\circ$ & $\circ$ & \textbf{New} & $\circ$ & $\circ$ \\
$\mathcal{P},\;\mathcal{P}_r$
  & ---
  & $\circ$ & $\circ$ & $\circ$ & \textbf{New} & $\circ$ \\
$\mathcal{Q}$
  & ---
  & $\circ$ & $\circ$ & $\circ$ & $\circ$ & \textbf{New} \\[4pt]

\multicolumn{7}{l}{\textit{Parameters}} \\
$\mu_{ij}$
  & $(i,j)$
  & $\circ$ & $+r$ & $\circ$ & $\times$ & $\circ$ \\
$ric_{ij}$
  & $(i,j)$
  & $\circ$ & $+r$ & $\circ$ & $\circ$ & $\circ$ \\
$bf_{ij},\;be_{ij},\;bi_{ij}$
  & $(i,j)$
  & $\circ$ & $+r$ & $\circ$ & $\times$ & $\circ$ \\
$c_{ij}$
  & $(i,j)$
  & $+m$ & $\circ$ & $\circ$ & $\circ$ & $\circ$ \\
$f_{ij}$
  & $(i,j)$
  & $+m$ & $\circ$ & $\circ$ & $\circ$ & $\circ$ \\
$v^{f},\;v^{e}$
  & scalar
  & $\circ$ & $+r$ & $\circ$ & $\circ$ & $\circ$ \\
$q^{\min},\;q^{\max}$
  & scalar
  & $+m$ & $\circ$ & $\circ$ & $\circ$ & $\times$ \\
$\omega_{ij},\;\delta_{ij}$
  & $(i,j)$
  & $+m$ & $\circ$ & $\circ$ & $\circ$ & $+q$ \\
$ssf,\;sse$
  & $(i,j)$
  & $\circ$ & $+r$ & $\circ$ & $\circ$ & $\circ$ \\
$s^{\text{RTI}},\;p^{\text{RTI}}$
  & scalar
  & $\circ$ & $+r$ & $\circ$ & $\circ$ & $\circ$ \\
$l_{ij},\;\alpha^f,\;\alpha^e$
  & $(i,j)$/scalar
  & $\circ$ & $\circ$ & $\circ$ & $\circ$ & $\circ$ \\
$d_{pij}$
  & ---
  & $\circ$ & $\circ$ & $\circ$ & \textbf{New} & $\circ$ \\
$\eta_{pr}$
  & ---
  & $\circ$ & $\circ$ & $\circ$ & \textbf{New} & $\circ$ \\
{$bm_{ijm}$}
  & {---}
  & {\textbf{New}} & {$\circ$}
  & {$\circ$}  & {$\circ$}
  & {$\circ$} \\
{$fhc_{i},\;vhc_{i}$}
  & {---}
  & {$\circ$}     & {$\circ$}
  & {\textbf{New}} & {$\circ$}
  & {$\circ$} \\[4pt]

\multicolumn{7}{l}{\textit{Binary variables}} \\
$X_{ij},\;XF_{ij},\;XI_{ij},\;XE_{ij}$
  & $(i,j)$
  & $+m$ & $+r$ & $\circ$ & $\circ$ & $+q$ \\[4pt]
$XH_{i}$
  & $(i)$
  & $\circ$ & $\circ$ & \textbf{New} & $\circ$ & $\circ$ \\[4pt]

\multicolumn{7}{l}{\textit{Flow variables}} \\
$NF_{ij},\;NI_{ij},\;NE_{ij}$
  & $(i,j)$
  & $+m$ & $+r$ & $\circ$ & $\circ$ & $+q$ \\
$q_{ij}$
  & $(i,j)$
  & $+m$ & $+r$ & $\circ$ & $\circ$ & $\times$ \\
$NS_{ij},\;NSF_{ij},\;NSI_{ij},\;NSE_{ij}$
  & $(i,j)$
  & $+m$ & $+r$ & $\circ$ & $\circ$ & $+q$ \\
$VH_{i}$
  & $(i)$
  & $\circ$ & $\circ$ & \textbf{New} & $\circ$ & $\circ$ \\[4pt]
$CH_{i}$
  & $(i)$
  & $\circ$ & $\circ$ & \textbf{New} & $\circ$ & $\circ$ \\[4pt]
$VS_{ij},\;VSF_{ij},\;VSI_{ij},\;VSE_{ij}$
  & $(i,j)$
  & $+m$ & $+r$ & $\circ$ & $\circ$ & $+q$ \\
$TNCF_{ij},\;TNCE_{ij},\;TNCI_{ij}$
  & $(i,j)$
  & $+m$ & $+r$ & $\circ$ & $\circ$ & $+q$ \\
$P$
  & scalar
  & $\circ$ & $+r$ & $\circ$ & $\circ$ & $\circ$ \\
$TC_{ij},\;TCS_{ij}$
  & $(i,j)$
  & $+m$ & $+r$ & $\circ$ & $\circ$ & $+q$ \\[4pt]
\end{longtable}

% ─────────────────────────────────────────────────────────────────────────────
%  NEW CONSTRAINTS PER EXTENSION
% ─────────────────────────────────────────────────────────────────────────────

\subsection{New Constraints Introduced by Each Extension}
\label{sec:ext-new-constraints}

% ── Extension M ──────────────────────────────────────────────────────────────
\paragraph{Extension M --- Transport Modes.}
{When multiple transport modes $m \in \mathcal{M}$ are
available on each route, the binary $X_{ijm}$ indicates whether mode $m$ is
used on route $(i,j)$.  The parameter $bm_{ijm} \in \{0,1\}$ encodes which
modes are physically admissible on each route.  Two rules govern mode
selection.}

\noindent\textit{{Admissibility.}}
{A mode may only be selected if it is admissible on that route:}
%
\begin{align}
  \overline{X}_{ijm} \leq bm_{ijm}
  && \forall\,(i,j)\in E,\;{m\in\mathcal{M}}
  \label{eq:ext-m-or}
\end{align}
%
{where $\overline{X}_{ijm}$ is the OR-aggregate binary
(Rule~3) indicating that mode $m$ carries at least one RTI type on route
$(i,j)$.}

\noindent\textit{{At-most-one mode per route.}}
Equation~\eqref{eq:route-enforced-by-parameters} is replaced by
\eqref{eq:at-most-one-m} to enforce {that at most one
transport mode is selected per route:}
\begin{align}
  \widehat{X}_{ij} \leq bf_{ij}+be_{ij}+bi_{ij}
  && {\forall\,(i,j)\in E}
  \label{eq:at-most-one-m}
\end{align}
{where $\widehat{X}_{ij}$ is the at-most-one aggregate binary
(Rule~4) equal to~1 if exactly one mode is active on route $(i,j)$.}

% ── Extension R ──────────────────────────────────────────────────────────────
\paragraph{Extension R --- RTI Types.}
{When multiple RTI types $r \in R_{ij}$ are available on
route $(i,j)$, the binary $XF_{ijr}$ indicates whether RTI type $r$ carries
full loads on that route.  Extension R replaces the single-type full-route
activation~\eqref{eq:full-activation} with a per-type rule.}

Each route commits to exactly one admissible RTI type (Rule~4); demand and
inlay flows are conditioned on the selected type.
Equation~\eqref{eq:full-activation} is replaced by~\eqref{eq:or-for-r} to
enforce only {\sout{different}}
{admissible} full RTI types {are sent on each active route:}
%
\begin{align}
  \overline{XF}_{ijr} = bf_{ijr}
  && \forall\,(i,j)\in E,\;{r \in R_{ij}}
  \label{eq:or-for-r}
\end{align}
%
{where $\overline{XF}_{ijr}$ is the OR-aggregate binary
(Rule~3) indicating that RTI type $r$ is used for full loads on route
$(i,j)$, and $bf_{ijr}=1$ if type $r$ may carry full loads on route
$(i,j)$.}

Equation~\eqref{eq:route-enforced-by-parameters} is replaced by
\eqref{eq:link-r-and-route}  {to ensure a route is only
activated when at least one admissible RTI type exists on it:}
\begin{align}
  \overline{X}_{ij} \leq \sum_{r\in R_{ij}} \bigl(bf_{ijr}+be_{ijr}+bi_{ijr}\bigr)
  && {\forall\,(i,j)\in E}
  \label{eq:link-r-and-route}
\end{align}
{where $\overline{X}_{ij}$ is the OR-aggregate binary
(Rule~3) equal to~1 if route $(i,j)$ carries any RTI type.}

% ── Extension H ──────────────────────────────────────────────────────────────
\paragraph{Extension H --- Hub Nodes.}
{Hub nodes $i \in H \subseteq V$ are intermediate
consolidation points that aggregate outgoing flows and incur a handling cost.
The parameter $fhc_i \geq 0$ is the fixed cost of activating hub $i$, and
$vhc_i \geq 0$ is the variable cost per unit volume handled at hub $i$.}

Hub nodes $i\in H$ are unconstrained. At every non-hub node
$i\in V\setminus H$ the incoming and outgoing empty-flow binaries are linked
to aggregated flows via Rule~3 and forced mutually exclusive. The index $r$
appears only when Ext.~R is also active:
%
\begin{align}
  \overline{XE}_{\cdot i} + \overline{XE}_{i \cdot} \leq 1
  && \forall\,i\in V\setminus H
  \label{eq:ext-h-excl}
\end{align}
%
{For every hub node $i \in H$, the following equations define
the hub activation binary $XH_i$, the total volume $VH_i$ flowing through
hub $i$, and the resulting hub cost $CH_i$:}
%
\begin{align}
  XH_i &= \overline{X}_{i\cdot}
  && {\forall\,i\in H}
  \label{eq:ext-h-XH} \\[4pt]
  VH_i &= V_{i\cdot}
  && {\forall\,i\in H}
  \label{eq:ext-h-VH} \\[4pt]
  CH_{i} &= fhc_{i}\cdot XH_i + vhc_i \cdot VH_i
  && {\forall\,i\in H}
  \label{eq:ext-h-CH}
\end{align}
%
{where $\overline{X}_{i\cdot} \coloneqq \sum_{j\in V}
\overline{X}_{ij}$ and $V_{i\cdot} \coloneqq \sum_{j\in V} V_{ij}$ are the
aggregated outgoing activation binary and total volume at node $i$.
The hub cost is added to the objective function~\eqref{eq:objective}:}
%
\begin{align}
  \min \quad
  \sum_{(i,j)\in E} TC_{ij}
  + p^{\text{RTI}} \cdot P
  + {\sum_{i\in H} CH_i}
  \label{eq:objective-H}
\end{align}

% ── Extension P ──────────────────────────────────────────────────────────────
\paragraph{Extension P --- Explicit Product Flows
  {(requires Extension R)}.}
%
{Extension P disaggregates the aggregate daily demand
$\mu_{ij}$ into product-level flows $NF_{ijpr}$, indexed by product
$p \in \mathcal{P}$ and RTI type $r \in R_{ij}$.  The parameter
$d_{ijp} \geq 0$ is the average daily demand of product $p$ on route $(i,j)$
in product units, and $\eta_{pr} > 0$ is the number of product-$p$ units
that fit in one RTI of type $r$.  Because the RTI-type index $r$ is
introduced by Extension R, Extension P cannot be used without it.}

{The average daily demand in RTI units is:}
%
\begin{align}
  \mu_{ijpr} = \frac{d_{ijp}}{\eta_{pr}}
  && {\forall\,(i,j)\in E,\;r\in R_{ij},\;p\in\mathcal{P}_r}
  \label{eq:ext-p-mu}
\end{align}
%
{The daily full-RTI flow for product $p$ using RTI type $r$
on route $(i,j)$ is:}
%
\begin{align}
  NF_{ijpr} = \mu_{ijpr} \cdot XF_{ijpr}
  && {\forall\,(i,j)\in E,\;r\in R_{ij},\;p\in\mathcal{P}_r}
  \label{eq:ext-p-NF}
\end{align}
%
{Equation~\eqref{eq:full-activation} is replaced by
\eqref{eq:ext-r-sel} to enforce that for each product $p$, at most one
admissible RTI type is selected per route:}
%
\begin{align}
  \widehat{XF}_{ijp} = bf_{ijp}
  && {\forall\,(i,j)\in E,\;p\in\mathcal{P}}
  \label{eq:ext-r-sel}
\end{align}
%
{where $\widehat{XF}_{ijp}$ is the at-most-one binary
(Rule~4) equal to~1 if product $p$ is assigned to exactly one RTI type on
route $(i,j)$, and $bf_{ijp}=1$ if at least one RTI type admissible for
product $p$ exists on route $(i,j)$.}

% ── Extension L ──────────────────────────────────────────────────────────────
\paragraph{Extension L --- Linearisation of products of two decision variables.}
